\documentclass[11pt]{combiArticle}
\usepackage{tikz,amsmath,amsfonts,latexsym}
\usepackage{amssymb}
\usetikzlibrary{shapes,arrows,calc}
\usepackage{calc}
\usepackage{parskip}
\usepackage{enumitem}
\setlist{itemsep=2pt, topsep=2pt, parsep=0pt}
\setlength{\emergencystretch}{2em}

\definecolor{myblue}{RGB}{0, 84, 159}

\lang{de}

\newcommand{\code}[1]{\texttt{#1}}
\newcommand{\taste}[1]{\fbox{\small\strut #1}}
\newcommand{\knopf}[1]{\textbf{#1}}

\begin{document}
\course{Schichtplan Manager}
\assistent{Linus Palm}
\contact{palm@combi.rwth-aachen.de}
\title{Betriebsanleitung}
\date{08. November 2026}
\maketitle

Der Schichtplan Manager plant den Dienstplan einer Station für einen ganzen Monat. Diese Demo-Version arbeitet
mit einer erfundenen \textbf{Musterstation}: neun Fachkräfte (drei davon in Teilzeit), eine Hilfskraft und zwei
Auszubildende. Sie nimmt Einschränkungen, Wünsche und die Mindestbesetzung auf und erstellt daraus in etwa
30 Sekunden einen geprüften Dienstplan. Die Anwendung läuft lokal auf dem MacBook und wird im Browser bedient.
Für die Berechnung braucht sie eine \textbf{Internetverbindung}, weil der Solver seine Lizenz online prüft.

%==========================================================================================================
\section{Programm starten}

\begin{enumerate}
  \item Den Mac einschalten und das Passwort eingeben (wird vor Ort mitgeteilt).
  \item Mit dem WLAN verbinden; ohne Internet kann kein Dienstplan berechnet werden.
  \item Warten, bis sich \textbf{Docker} und \textbf{Visual Studio Code} von selbst geöffnet haben.
  \item In \textbf{Visual Studio Code} oben im Menü \emph{Terminal $\to$ Neues Terminal} wählen. Unten öffnet sich ein
        Terminal.
  \item Dort \code{git pull} eingeben und mit \taste{Enter} bestätigen, um die aktuelle Version des Schichtplaners zu laden.
  \item Dann \code{just run} eingeben und mit \taste{Enter} bestätigen, um ihn zu starten.
  \item Warten, bis alles aufgebaut ist: Die Meldungen zu \code{api} und \code{webapp} enden grün mit
        \emph{Healthy}, danach erscheint wieder die Eingabezeile. Beim ersten Start kann das einige Minuten
        dauern. Warnungen ganz am Anfang, etwa dass \code{pnpm} oder \code{node} nicht gefunden wurde, sind
        unbedenklich.
  \item Mit \taste{cmd} + \taste{Leertaste} die Suche öffnen, \code{safari} eingeben und mit \taste{Enter} den
        Browser öffnen.
  \item Entweder \code{http://localhost:3000} aufrufen oder in den Lesezeichen
        \emph{Übersicht-Schichtplan Manager} anwählen. Die Übersicht des Schichtplan Managers erscheint.
\end{enumerate}

\section{Programm beenden}

Im selben Terminal in Visual Studio Code \code{just stop} eingeben und mit \taste{Enter} bestätigen. Danach
kann der Mac wie gewohnt ausgeschaltet werden.

\textbf{Achtung:} Alles, was eingetragen oder berechnet wurde (Einschränkungen, Wünsche, Mindestbesetzung und
Dienstpläne), wird nur gehalten, solange das Programm läuft. Nach einem Neustart beginnt die Musterstation wieder
im Ausgangszustand. Einen Dienstplan, den man behalten möchte, vorher speichern (Abschnitt~\ref{sec:pruefen}).

%==========================================================================================================
\section{Aufbau der Oberfläche}

\begin{itemize}
  \item \textbf{Seitenleiste} links: \emph{Übersicht}, unter \emph{Planungsdaten} die Seiten
        \emph{Mitarbeiter}, \emph{Verfügbarkeit} und \emph{Mindestbesetzung}, unter \emph{Dienstplan} die Seiten
        \emph{Erstellen} und \emph{Prüfen}. Jede Seite ist jederzeit erreichbar.
  \item \textbf{Monat und Jahr} neben dem Seitentitel. Alle Seiten arbeiten mit diesem Monat; er bleibt beim
        Seitenwechsel erhalten. Jeder Monat der Jahre 2026 und 2027 kann geplant werden, immer als ganzer Monat.
        Beim Öffnen ist der nächste Monat gewählt.
  \item \textbf{Übersicht}: eine Karte je Seite in der üblichen Reihenfolge der Planung. Der Pfeil vor einem
        Seitentitel führt zur Übersicht zurück.
\end{itemize}

Die übliche Reihenfolge: Mitarbeiter ansehen $\to$ Verfügbarkeit und Mindestbesetzung pflegen $\to$
Dienstplan erstellen $\to$ Dienstplan prüfen.

%==========================================================================================================
\section{Mitarbeiter}

Zeigt die zwölf Mitarbeiter der Musterstation, nur zum Ansehen.
\begin{itemize}
  \item Suche nach Name, Personalnummer oder Qualifikation.
  \item Eine Zeile aufklappen zeigt die \textbf{Zuordnung} mit der Qualifikation, das \textbf{Monatskonto}
        (Soll und Gutschriften mit Datum) und die \textbf{Verfügbarkeit} im Monat.
\end{itemize}
Das Soll ergibt sich aus der Arbeitszeit (Vollzeit 38,5 Stunden pro Woche, Teilzeit anteilig) und den Werktagen
des Monats. Urlaub und Fortbildung werden mit dem Tagessoll gutgeschrieben.

%==========================================================================================================
\section{Verfügbarkeit}

Hier werden für einen Mitarbeiter (Auswahl oben auf der Seite) \textbf{Einschränkungen} und \textbf{Wünsche}
im Monatskalender gepflegt. Der Kalender zeigt je Tag zwei Dinge:
\begin{itemize}
  \item \textbf{Einschränkung}, rot: \emph{Nicht verfügbar}, \emph{Urlaub}, \emph{Fortbildung}, \emph{Frei}
        oder \emph{Nur bestimmte Schichten} mit den erlaubten Schichten, dazu ein optionaler Grund.
        Einschränkungen werden bei der Planung \textbf{immer} eingehalten.
  \item \textbf{Wunsch}, grün: freier Tag, bestimmte Schicht frei, Wunschtag oder Wunschschicht. Wünsche werden
        so weit wie möglich erfüllt, aber nicht garantiert.
\end{itemize}
Einen Tag anklicken, rechts Einschränkung oder Wunsch wählen und \knopf{Einschränkung speichern} bzw.
\knopf{Wunsch speichern} drücken; \knopf{Einschränkung entfernen} bzw.\ \knopf{Wunsch entfernen} löscht den
Eintrag. Pro Tag gibt es höchstens eine
Einschränkung und einen Wunsch. \textbf{Mehrere Tage auf einmal:} \taste{Umschalt}-Klick wählt einen Zeitraum,
\taste{ctrl}-Klick fügt einzelne Tage hinzu oder entfernt sie; das Speichern gilt dann für alle gewählten Tage. Die Hinweise über dem Kalender zeigen die Tasten.

Jeder Monat beginnt mit einigen Beispieleinträgen: einer Urlaubswoche, einem Fortbildungstag und zwei Wünschen.
Sie lassen sich wie jeder andere Eintrag ändern oder entfernen.

%==========================================================================================================
\section{Mindestbesetzung}

Legt fest, wie viele Mitarbeiter je Tag, Schicht und Qualifikation mindestens im Dienst sein müssen. Jeder Monat
beginnt mit einer Vorgabe, die zum Team passt.
\begin{itemize}
  \item Reiter wählen die Qualifikation: \emph{Fachkraft}, \emph{Hilfskraft}, \emph{Azubi}.
  \item Das Raster hat eine Zeile je Tag und eine Spalte je Schicht (F, Z, S, N). Wochenenden und Feiertage
        (NRW) sind hinterlegt. Jedes Feld bietet 0 bis 5 zur Auswahl; 0 heißt: niemand nötig.
  \item \textbf{Wochenmuster anwenden} füllt den Monat schnell: Werte für Montag bis Sonntag und für Feiertage
        wählen, \knopf{Vorschau} zeigt die geänderten Tage, \knopf{Übernehmen} trägt sie ins Raster ein.
  \item Geänderte Zellen sind orange umrandet. \knopf{Speichern} sichert den Monat, \knopf{Zurücksetzen} stellt
        den gespeicherten Stand wieder her.
\end{itemize}
\textbf{Wichtig:} Ungespeicherte Änderungen gehen beim Verlassen der Seite verloren. Mehr Bedarf, als das Team
leisten kann, bleibt im Dienstplan als Lücke offen.

%==========================================================================================================
\section{Dienstplan erstellen}

\begin{enumerate}
  \item Den Monat oben wählen. Die Karte nennt, was geplant wird, und listet unter \emph{Berücksichtigt}, welche
        Daten und Regeln in die Planung eingehen.
  \item \knopf{Dienstplan erstellen} drücken.
\end{enumerate}
\textbf{Vormonat (optional):} Unter \emph{Vormonat} lässt sich der gespeicherte Dienstplan des Vormonats laden,
z.\,B.\ \code{Dienstplan November 2026 - Version 2.json} für den Dezember. Dann gelten Nachtdienste, Ruhezeiten,
Arbeitstage am Stück und Wochenenden auch über den Monatswechsel. Angenommen wird nur ein unveränderter Plan
genau des Vormonats, der alle Regeln einhält; sonst erscheint eine Meldung. \knopf{Entfernen} nimmt ihn wieder
heraus. Ohne Vormonat gelten die Tage vor dem Monat als frei.

Die Berechnung dauert etwa 30 Sekunden; ein Balken zeigt die verbleibende Zeit. Die Seite darf währenddessen
verlassen werden. Unter \emph{Letzte Berechnung} steht dann das Ergebnis, mit einer Überschrift, die sagt, was zu
tun ist:
\begin{itemize}
  \item \emph{Dienstplan erstellt, Regeln eingehalten}: Mit \knopf{Dienstplan prüfen} weiter zur Prüfung.
  \item \emph{Berechnung nicht möglich}: Der Solver kann seine Lizenz nicht prüfen, meist fehlt die Internetverbindung. WLAN prüfen und erneut starten.
  \item \emph{\dots\ mit Regelverstößen} oder \emph{\dots\ unvollständig geprüft}: Der Kasten \emph{Verstöße}
        bzw.\ \emph{Hinweise} nennt die Ursache.
\end{itemize}
Darunter stehen Zeit, Ergebnis der Prüfung, Anzahl der Dienste und der \textbf{Lücken}, also Pflichtstellen, die
mangels Personal offen bleiben. Ein Plan mit Lücken kann trotzdem gültig sein. Es läuft immer nur eine Berechnung
gleichzeitig. Wie der Plan entsteht, beschreibt die \emph{Modell -- Kurzfassung}.

%==========================================================================================================
\section{Dienstplan prüfen}\label{sec:pruefen}

Zeigt den zuletzt erstellten oder geladenen Dienstplan. Von oben nach unten:
\begin{itemize}
  \item \textbf{Zusammenfassung} mit der Statuszeile. Nur bei \emph{Regeln eingehalten} ist der Plan
        verwendbar. Bei \emph{Regelverstöße} oder \emph{Unvollständig geprüft} listet ein
        Kasten jeden Befund mit Regel, Mitarbeiter und Datum.
  \item \textbf{Lücken}: je Tag, Schicht und Qualifikation, wie viele Mitarbeiter fehlen. Für diese Schichten
        Gastpersonal anfragen.
  \item \textbf{Kennzahlen}: der Plan auf einen Blick: Lücken, Monatssaldo, erfüllte Wünsche und Wechsel vom
        Spät- in den Frühdienst.
  \item \textbf{Wünsche} (aufklappbar): jeder Wunsch mit \emph{erfüllt}, \emph{nicht erfüllt} oder \emph{nicht
        erfüllbar} (z.\,B. wegen einer Einschränkung am selben Tag).
  \item \textbf{Dienstplan}: je Mitarbeiter und Tag die Schicht in ihrer Farbe: F gelb, Z türkis, S orange,
        N schwarz. Zeigt man auf einen Dienst, erscheinen Uhrzeit und Qualifikation. Kürzel zeigen
        Einschränkungen, rot
        hervorgehoben \code{U} Urlaub, \code{FB} Fortbildung und z.\,B.\ \code{nur F,S} für nur bestimmte
        Schichten. Wochenenden sind rosa hinterlegt.
        Suche nach Name und \knopf{Vollbild} helfen bei großen Plänen.
  \item \textbf{Besetzung} unter dem Plan: je Schicht und Tag \emph{Ist/Soll}, z.\,B.\ \code{3/2}.
        Rot hinterlegte Zellen sind Lücken.
  \item \textbf{Monatskonten} (aufklappbar): je Mitarbeiter Soll, Gutschriften, geplante Stunden und Saldo.
\end{itemize}
Die Knöpfe in der Zusammenfassung:
\begin{itemize}
  \item \knopf{Speichern}: speichert den Plan als eine Datei im Ordner \emph{Downloads}, benannt nach Monat, Jahr
        und Version, z.\,B.\ \code{Dienstplan Dezember 2026 - Version 2.json}. Jeder neu erstellte Plan eines
        Monats bekommt die nächste Version, so lassen sich mehrere Pläne auseinanderhalten.
  \item \knopf{Laden}: öffnet einen gespeicherten Plan wieder; dazu die Datei wählen. So lassen sich mehrere
        Pläne nacheinander ansehen und vergleichen. Die vier vorbereiteten Pläne liegen als
        \code{Dienstplan Januar 2026 - Version 1} bis \code{4} in den Ordnern \emph{Plan 1} bis \emph{Plan 4}.
\end{itemize}
Die Version steht auch in der Zusammenfassung über dem Plan. Nach einem Neustart des Programms beginnt die
Zählung wieder bei 1, sobald kein gespeicherter Plan des Monats geladen wurde.

Gehört der Plan zu einem anderen Monat als dem oben gewählten, wechselt \knopf{Zu \dots\ wechseln} zu seinem
Monat.

%==========================================================================================================
\section{Wenn etwas nicht klappt}

\begin{itemize}
  \item \textbf{\code{just run} bricht ab:} Prüfen, ob Docker geöffnet ist und läuft, dann \code{just run}
        erneut eingeben.
  \item \textbf{Die Seite lädt nicht:} Warten, bis im Terminal alles grün ist, dann die Seite neu laden.
  \item \textbf{\emph{Berechnung nicht möglich}:} Internetverbindung prüfen, dann erneut
        \knopf{Dienstplan erstellen} drücken.
  \item \textbf{\emph{Die Anwendung antwortet nicht}:} kurz warten und die Seite neu laden.
  \item \textbf{Ganz neu starten:} \code{just stop}, danach \code{just run}.
\end{itemize}

\end{document}
